\section{Discussion and Limitations}
\label{sec:limitations}

A single temporal execution rule is not behaviorally neutral across the
components of the evaluated ACT actions. The preregistered same-target contrast
shows a large arm--gripper asymmetry, the development sweep preserves its
ordering over 0.10--1.60~s, and the within-arm comparison separates translation
from rotation. The practical consequence is narrower: longer gripper
commitment improves two schedules when arm cadence is fixed, but coherent H16
remains the strongest overall frozen operating point.

Same-target alignment preserves the physical step for which every command was
predicted, but it couples two temporal quantities. An older source observation
necessarily contributes a farther-lookahead prediction because $k=d$. The
diagnostic therefore estimates their joint closed-loop effect and cannot say
whether source age, prediction horizon, or their combination produces the
behavioral difference. This limitation is part of the intervention definition,
not merely a statistical qualification.

The age sweep and translation--rotation comparison characterize an Object
development cohort rather than the primary confirmatory cohort. Their blocks
are exactly matched across the conditions being compared, but historical
anchors and later reviewer-directed probes were generated at different study
stages. The separate Spatial completion preserves the main asymmetry's sign,
yet it was also conducted after the primary study and was not pooled with the
140-block result.

Same-target probes are diagnostic interventions, not candidate deployment
schedules. They query the policy at every executed step even when an older
prediction supplies part of the action. The dense-query factorial shows that
discarded queries alone cannot explain one conditional H16 effect under the
tested deterministic evaluator. Track A then supplies the deployment-facing
comparison at fixed arm cadences and nearly matched policy-query rates. Query
rate is not a measurement of equal latency, total query count, or equal
compute.

The Track-A gains are heterogeneous across suites. Both component-resolved
contrasts are largest in LIBERO-10, while Goal and Spatial operate near the H16
ceiling. Because suite identity, baseline difficulty, and task semantics vary
together, their individual roles cannot be identified. This heterogeneity and
the 357/450 H16 result rule out a claim that component-resolved scheduling is a
globally superior executor in the evaluated static benchmark.

The tested diagnostics do not identify a positive mechanism. Same-target
dispersion predicts the wrong within-arm ordering, gripper persistence is
incomplete, the occupancy moderator is null, forecastability lacks a
prespecified discriminative curve pattern, and command discontinuity cannot
separate the competing explanations. These results narrow the explanation
space without resolving why translation, rotation, and gripper commands respond
differently.

Finally, the central behavioral evidence is specific to the evaluated ACT
policies, LIBERO tasks, checkpoints, and control interface. SmolVLA
~\cite{shukor2025smolvla} cannot
provide a physically matched replication because its training and chunk
physical timebase cannot be identified from the available provenance. A
planned SmolVLA evaluation was not run because its prespecified runtime
eligibility window had expired. We therefore make no cross-policy or
physical-robot claim. Broader action spaces, stochastic evaluators, and real
robots remain outside the evidence reported here.
